\documentclass[./eval.tex]{subfiles}
\newcommand{\head}{\noindent Máquinas Eléctricas II |Recupertorio Evaluación| Apellidos y Nombres.
\rule{10cm}{0.1mm}}
\newcommand{\separador}{\noindent\rule{\paperwidth}{0.1mm}\newline}
\begin{document}
\head

 \begin{enumerate}
 \item Un conductor se desplaza a una velocidad lineal de $v=4m/s$ dentro de un campo magnético fijo de $B=10T$ de inducción. Determinar el valor de la f.e.m. inducida en el mismo si posee una longitud de $l=0,5m$.
    \item A que velocidad se deberá mover un conductor de longitud $l=0.4m$ para que dentro de una inducción de campo magnético de $B=2.5T$ produzca una f.e.m de $e=5V$.
    \item ¿Cuál la f.e.m de auto-inducción si el flujo magnético $\Phi$ que existe en el campo magnético producido por una bobina de una electroválvula si tiene esta un núcleo cilíndrico de diámetro $d=0,04m$ y la inducción magnética en la misma crece de $B=0T$ a $B=1.5T$ en $\Delta t=0,3s$?
    \item Calcular el valor de la f.e.m de auto-inducción que se desarrollará una boina con coeficiente de auto-inducción de 50 milihernios si se aplica una corriente que crece desde cero hasta $I=10A$ en un tiempo de $\Delta t=0,1 seg.$
    \separador
\end{enumerate}
\end{document} 